% ------------------------------------------------------------------------------
% 1 Introduction
% ------------------------------------------------------------------------------
\section{Introduction}\label{sec:intro}

Colorectal cancer is one of the most common malignancies worldwide, and for
patients undergoing potentially curative resection the question of whether
systemic adjuvant chemotherapy should follow surgery remained open for much
of the twentieth century. The North Central Cancer Treatment Group (NCCTG)
program addressed this question directly: building on an earlier randomized
comparison of observation against levamisole alone~\cite{laurie1989}, the
investigators added a third arm combining the immunomodulator levamisole
with the antimetabolite 5-fluorouracil (5-FU)~\cite{moertel1990}. The
combination reduced recurrence and mortality substantially and became a
standard of care for resected stage B/C colon carcinoma. The individual
patient data from this program were subsequently contributed to the
\texttt{survival} package for R~\cite{therneau2026}, where they now serve as
one of the canonical datasets for teaching survival analysis; the version
analysed here contains 929 patients from a single study.

Re-analyzing an old trial with new methods is more than an academic
exercise, and five considerations motivated each analytical module of this
project. First, endpoint definitions and analysis conventions have evolved:
modern adjuvant colorectal trials define disease-free survival (DFS) as the
time to recurrence \emph{or} death, whereas the published coding of this
dataset treats death without recurrence as censoring; a rigorous
re-analysis must therefore address both views and, importantly, the
competing-risks consequences of the difference. Second, missing-data
handling, covariate-adjusted estimation, and absolute-benefit reporting
have all been refined since the original publications, and regulatory
guidance now actively encourages pre-specified covariate adjustment in the
primary analysis of randomized trials~\cite{fda2023,kahan2014}. Third, the
trial-design landscape has changed: group-sequential designs with alpha-
and beta-spending functions are now the default in oncology, and it is
natural to ask whether this trial---run with a fixed-sample
mindset---would have stopped early under a modern monitoring
plan~\cite{lan-demets1983,o1979}. Fourth, machine-learning prediction
models and Bayesian machinery offer perspectives on prognosis and evidence
accumulation that complement the classical
toolkit~\cite{harrell2015,gelman2013}. Finally, the modern estimand
framework and recent methods for heterogeneous treatment effects invite a
re-examination of \emph{for whom} the benefit is
largest~\cite{hernan2020}.

Accordingly, this project integrates all of these perspectives into a
single coherent pipeline with six objectives:

\begin{enumerate}[leftmargin=1.7em, itemsep=1.5pt, topsep=3pt]
  \item \textbf{Re-estimate the treatment effects} on DFS and overall
        survival with covariate-adjusted Cox models fitted on multiply
        imputed data, explicit proportional-hazards diagnostics,
        competing-risk sensitivity analyses, and bootstrap and
        restricted-mean-survival-time absolute-benefit measures
        (Sections~\ref{sec:methods-freq}~and~\ref{sec:res-tx});
  \item \textbf{Redesign the trial} as it would be designed today---a
        formal sample-size calculation and a three-look group-sequential
        design---and re-enact interim monitoring on the real event times
        (Sections~\ref{sec:methods-design}~and~\ref{sec:res-design});
  \item \textbf{Quantify operating characteristics by simulation},
        covering the efficiency of covariate adjustment, the type-I error
        and power of the sequential design, and confidence-interval
        coverage under non-proportional hazards
        (Sections~\ref{sec:methods-sim}~and~\ref{sec:res-sim});
  \item \textbf{Build and honestly validate prognostic survival models}
        spanning classical regression, penalization, and ensembles of
        trees (Sections~\ref{sec:methods-ml}~and~\ref{sec:res-ml});
  \item \textbf{Conduct a Bayesian re-analysis} with prior sensitivity,
        event-driven sequential monitoring, posterior predictive checks,
        and baseline-model criticism
        (Sections~\ref{sec:methods-bayes}~and~\ref{sec:res-bayes});
  \item \textbf{Characterise heterogeneous treatment effects} on the
        absolute five-year risk using jackknife pseudo-observations,
        cross-fitted G-computation, and doubly robust AIPW estimation
        (Sections~\ref{sec:methods-cate}~and~\ref{sec:res-cate}).
\end{enumerate}

The report is organized to mirror the repository structure: Section
\ref{sec:data} describes the data and endpoint construction; Section
\ref{sec:methods} details all statistical methods; Sections
\ref{sec:res-tx}--\ref{sec:res-cate} present results in five thematic
blocks; Section~\ref{sec:discussion} discusses findings, strengths, and
limitations. All analyses are scripted in R and Python with fixed random
seeds, and Section~\ref{sec:software} records the exact software
environment required for byte-level reproduction.

% ------------------------------------------------------------------------------
% 2 Data
% ------------------------------------------------------------------------------
\section{Data}\label{sec:data}

\subsection{Source and patients}

The analysis dataset derives from the \texttt{colon} data frame
distributed with the R \texttt{survival} package~\cite{therneau2026},
obtained as a CSV file from the Rdatasets mirror and stored in the
repository under \texttt{data/raw/}. The underlying trial was conducted by
the NCCTG and reported by Laurie et al.~\cite{laurie1989} and Moertel et
al.~\cite{moertel1990}. The file stores \emph{two records per patient}: one
for time-to-recurrence and one for time-to-death, with baseline covariates
duplicated across the pair. After verifying covariate consistency within
each patient (all 929 patients passed), the two records were pivoted into a
single patient-level analysis file of 929 rows.

Patients were randomized 1:1:1 within a single study to observation
(\texttt{Obs}, $n=315$), levamisole alone (\texttt{Lev}, $n=310$), or
levamisole plus 5-FU (\texttt{Lev+5FU}, $n=304$) following surgical
resection. Baseline covariates available for adjustment and prediction
comprise: age (mean 59.8 years), sex (47.9\% female), obstruction of the
colon by tumour (19.4\%), perforation (2.9\%), adherence to nearby organs
(14.5\%), number of positive lymph nodes (median 2, IQR 1--5), tumour
differentiation (well/moderate/poor), extent of local spread (submucosa
through contiguous structures), and the surgery-to-registration interval
category (short/long).

\subsection{Endpoint construction}\label{sec:endpoints}

The published coding defines the \texttt{etype = 1} record as time to
tumour recurrence, with death before recurrence treated as censoring. In
the present dataset this coding produces 468 recurrences and, notably, 38
deaths that occur without a preceding recurrence flag at the same visit
date. Because contemporary adjuvant colorectal trials define DFS as
recurrence \emph{or} death from any cause, the data-curation module
constructed three complementary endpoint views:

\begin{enumerate}[leftmargin=1.7em, itemsep=1.5pt, topsep=3pt]
  \item \textbf{Composite DFS} (primary endpoint): time to first of
        recurrence or death, whichever occurs first; 506 events (54.5\%).
  \item \textbf{Overall survival} (secondary endpoint): time to death from
        any cause; 452 events (48.7\%).
  \item \textbf{Recurrence with competing events} (sensitivity analyses):
        cause-coded status with 468 recurrences (cause 1), 38 competing
        deaths (cause 2), and 423 censored observations, supporting
        Aalen--Johansen estimation and Fine--Gray subdistribution
        modelling.
\end{enumerate}

Cross-endpoint consistency checks confirmed that the recurrence time never
exceeds the death time and that all 38 competing deaths occur exactly at
their death dates. Follow-up is substantial: the median follow-up among
censored patients, estimated by reverse Kaplan--Meier, is 6.4 years, and
the maximum observed time is 9.1 years.

\subsection{Missing data and imputation}

Two adjustment covariates contain missing values: the number of positive
lymph nodes (18 patients, 1.9\%) and tumour differentiation (23 patients,
2.5\%); 95.6\% of patients are complete cases. Following White and
Royston's recommendation for survival outcomes~\cite{white2009}, missing
covariates were imputed by fully conditional specification (chained
equations) using the \texttt{mice} package~\cite{mice2011} with $m=20$
imputations, including the Nelson--Aalen cumulative hazard estimate and
the event indicator of the relevant endpoint in each imputation model so
that the imputation is compatible with the subsequent Cox analysis. A
plausibility check comparing observed and imputed distributions showed
modest and clinically reasonable differences (mean imputed nodes 3.3 vs.\
observed 3.7), and pooled estimates were
contrasted against complete-case results throughout
(Section~\ref{sec:res-tx}).

\subsection{Baseline balance}

Randomization produced well-balanced arms. Across eleven baseline
variables, the only nominal imbalance was sex (53.6\% female in
\texttt{Lev+5FU} vs.\ 47.3\% in \texttt{Obs}; three-arm global exact
test $p=0.027$, pairwise $p=0.13$), which is unremarkable in the context
of eleven simultaneous tests and does not survive multiplicity
correction. Pairwise standardized mean differences
(Figure~\ref{fig:smd}) never exceeded $|0.216|$, and the overall picture
is consistent with adequate randomization concealment.

\begin{figure}[htbp]
  \centering
  \includegraphics[width=0.95\textwidth]{fig_smd_loveplot.pdf}
  \caption{Baseline covariate balance across randomized arms, measured by
  absolute standardized mean differences (SMD). Dashed lines mark the
  conventional $|\mathrm{SMD}|=0.10$ threshold. All pairwise comparisons
  between arms are shown.}
  \label{fig:smd}
\end{figure}
